%% ==========================================================================
%%  FINAL REPORT --- SECTION 1, INTRODUCTION
%%  Self-contained: the survey's introduction.tex describes the survey and is
%%  left untouched, since it is the source of a submitted document.
%% ==========================================================================
\section{Introduction}
\label{sec:intro}

A recommender system is built from the most revealing data its users produce:
what they watched, read and bought. That this data is identifying, and not
merely sensitive, was shown by Narayanan and Shmatikov~\cite{narayanan2008}, who
matched records in the anonymised Netflix Prize dataset to public profiles and
report ``uncovering their apparent political preferences and other potentially
sensitive information''. A recommender that its operator cannot read therefore
has to protect more than one thing. It has to protect the ratings the model is
trained on, and it has to protect what the user does with the recommendations
afterwards, because fetching a recommended item says which item it was.

\subsection{Two papers, two halves}
\label{sec:intro:gap}

The course handed us two systems that each address one of these. \nudge{}
trains a matrix factorisation of the rating matrix across three servers, at most
one of them compromised, under 2-of-3 replicated secret sharing. Its central
observation is that a matrix--vector product under this sharing costs
communication proportional to the output vector rather than to the matrix, so an
algorithm written as a chain of such products with a few non-linear steps
between them is cheap. Power iteration has that shape. \nudgeplain{} publishes
the resulting item embeddings $\itemembed$ in the clear, by design, and computes
each user's scores in shared form. It then stops, and says so: the system
``relies on other means (e.g., Apple's private relay, Tor, or cryptographic
private information retrieval) to let users fetch data items in a private
way''~\cite{nudge2026}.

\pirsona{} addresses the fetch. Users retrieve items by multi-server private
information retrieval (PIR), and the servers harvest a secret-shared record of
what each user consumed directly from the PIR queries, so the next round of
training needs no separate upload. Its training core is a ``carefully optimized
4PC Boolean matrix factorization'', and its authors are explicit that they
``opted for the most performant primitives available (at the expense of rather
strong non-collusion assumptions)''~\cite{pirsona2021}.

Neither paper builds both halves well, and neither measures what happens when
they meet. That meeting point is where this project sits. \nudgeplain{}'s public
$\itemembed$ is a complete latent-factor model of the catalogue, known to
everyone, so a fetch observed in the clear is not an isolated fact about a user.
It is a projection of that user's private embedding onto a basis the observer
already holds. Section~\ref{sec:leak} measures how quickly those projections
give the user away.

\subsection{What we built and what we found}
\label{sec:intro:contrib}

\sys{} is the composition: \nudgeplain{}'s training core, followed by
two-server PIR built from a distributed point function (DPF), with
\pirsonaplain{}'s harvest connecting delivery back to training. Everything
cryptographic is written from scratch in C++17, including the DPF, which serves
both as the read layer and as the comparison gate inside training. Our
contributions are these.

\begin{enumerate}
  \item \textbf{A working composition.} Private training runs end to end on
    MovieLens-100K, the trained model is served from three separate server
    processes over TCP, and each recommended record is fetched by DPF-PIR and
    checked byte for byte against a cleartext lookup
    (Sections~\ref{sec:design} and~\ref{sec:impl}).

  \item \textbf{The cost of the composition, measured half by half.} Composing
    costs exactly the sum of the halves, which the benchmark asserts to the
    byte. \nudgeplain{}'s Theorem~4.2 is measured rather than quoted:
    quadrupling the input matrix raises traffic by 4.4\% and leaves the round
    count unchanged. Truncation accounts for 68.6 to 74.4\% of all rounds, and
    the FSS comparison gate for 3.1\%. A normaliser that reveals nothing costs
    between 0.0018 and 0.0041 of nDCG@20 and 5.3 times the rounds, and it
    cannot run in a 64-bit ring at all (Section~\ref{sec:eval}).

  \item \textbf{Why the fetch must be private, and why cheaper defences fail.}
    Ten observed fetches recover a user's private embedding to a cosine of
    0.84, well above a popularity control. Decoy traffic at five times the
    real volume leaves the attack far above that control, and \nudgeplain{}'s
    own differential-privacy mechanism makes the attack stronger while
    removing 71\% of recommendation quality. We reconcile that figure with
    \nudgeplain{}'s own 17\% (Section~\ref{sec:leak}).

  \item \textbf{A malicious-client hole in the composition, and its partial
    fix.} The harvest path accepted a vote of any weight. We measure what that
    buys an attacker, close weight inflation with a one-round check, show by
    construction what the check does not close, and measure the
    denial-of-service amplification that full-domain evaluation invites
    (Section~\ref{sec:leak:malicious}).

  \item \textbf{A stated security argument.} A threat model with six adversary
    classes, a leakage profile, and a real/ideal simulation sketch that walks
    every message. The sketch shows that one of our two normalisers is not
    simulable against the intended functionality, and we say which numbers
    came from it (Section~\ref{sec:threat}).
\end{enumerate}

We also report results that go against us, because a reader deciding whether to
build this needs them. On a wide-area link, downloading the entire catalogue
beats our private read by a factor of 4.6, since delivery costs one round trip
per recommendation and we did not batch them. The spec-faithful normaliser loses
most of its quality at a single power iteration. And one bug, a silently
clamping comparison range, would have been reported as ``privacy costs 24\% of
recommendation quality'' had we measured only cost.

\paragraph{What this report does not claim.} Security is semi-honest, and two
colluding servers learn everything. The evaluation runs on MovieLens-100K, not
on MovieLens-1M or Netflix scale. Wide-area profiles are emulated at the channel
level, not with kernel network emulation. The security argument is a sketch.
Section~\ref{sec:limits} states each of these in full.

\subsection{Organisation}
\label{sec:intro:org}

Section~\ref{sec:threat} fixes the threat model and the security argument.
Section~\ref{sec:background} covers the primitives, and
Section~\ref{sec:design} the system design. Section~\ref{sec:impl} describes the
implementation and how it was verified, and Section~\ref{sec:eval} evaluates its
cost and quality. Section~\ref{sec:leak} measures what the composition leaks and
what a malicious client can do. Section~\ref{sec:related} places the work,
Section~\ref{sec:limits} lists its limitations, and
Section~\ref{sec:conclusion} concludes. Appendix~\ref{app:protocols} gives the
protocols as pseudocode, and Appendix~\ref{app:reproduce} maps every figure to
the command and the data file that produce it.
